\section{Standard and costandard modules}\label{sec:standards}

We study the standard and costandard modules associated with the
quasi-hereditary structure induced by the Reedy algebra
$(\Lambda_P,\Lambda_P^+,\Lambda_P^-)$ constructed in the previous section.
We give a combinatorial description of the standard modules and of morphisms among them 
(\Cref{prop:standard-concrete-model,prop:standard-hom-formula}).
The corresponding descriptions for costandard modules are obtained by
duality.
We also use the filtration $J_\bullet$ of $\Lambda_P$ in
\eqref{eq:interval-cardinality-heredity-chain} to construct
filtrations of the indecomposable projective and injective modules
with standard and costandard subquotients, respectively, and to identify the successive quotients $J_t/J_{t-1}$ as bimodules (\Cref{sec:standard-cardinality-filtrations}).

\subsection{Standard modules}\label{sec:standard-modules}

For each
$S\in\Int(P)$, recall the canonical exact sequence
\begin{equation}
0
\longrightarrow
\KK(S)
\longrightarrow
\PP(S)
\xrightarrow{\pi_S}
\DD(S)
\longrightarrow
0.
\label{eq:interval-standard-defining-sequence}
\end{equation}
We write $[f]:=\pi_S(f)$ for $f\in\PP(S)$.


\subsubsection{Basis descriptions and algebra actions}
\label{sec:standard-concrete}
Let $S\in \Int(P)$. 
We first describe the kernel $\KK(S)$ in terms of the
admissible-component basis. 
\begin{lemma}
\label{lem:standard-kernel-admissible-basis}
For $R\in\Int(P)$, 
\begin{equation}
e_R\KK(S)
=
\bigoplus_{\substack{U\in\Adm(R,S)\\ U\subsetneq S}}
\kk\rho_{R,S}^{U}.
\label{eq:standard-kernel-admissible-basis}
\end{equation}
Moreover, the kernel is generated by the canonical maps associated with proper relative upsets of $S$:
\begin{equation}
\KK(S)
=
\sum_{\substack{C\in\cU(S)\\ C\subsetneq S}}
\Image (j_{C,S})_*.
\label{eq:standard-kernel-generators}
\end{equation}
\end{lemma}



\begin{proof}
By \eqref{eq:Reedy-standard-kernel} and $\deg(T)=|T|$, the space
$e_R\KK(S)$ is spanned by morphisms $\I_R\to\I_S$ that factor through
some $\I_T$ with $|T|<|S|$.  For such a factorization, it suffices to
consider
$\rho_{T,S}^{B}\circ\rho_{R,T}^{A}$ with
$A\in\Adm(R,T)$ and $B\in\Adm(T,S)$.  By
\Cref{prop:admissible-basis-composition},
\[
\rho_{T,S}^{B}\circ\rho_{R,T}^{A}
=
\sum_{W\in\piZero(A\cap B)}\rho_{R,S}^{W}.
\]
Each $W$ belongs to $\Adm(R,S)$ and satisfies
$W\subseteq T$, hence $W\subsetneq S$.  This gives one inclusion in
\eqref{eq:standard-kernel-admissible-basis}.  Conversely, if
$U\in\Adm(R,S)$ and $U\subsetneq S$, then
the morphism $\rho_{R,S}^{U}$
factors through $\I_U$, with $|U|<|S|$.  Thus
$\rho_{R,S}^{U}\in e_R\KK(S)$, proving
\eqref{eq:standard-kernel-admissible-basis}.

For every proper $C\in\cU(S)$, the image of $(j_{C,S})_*$ is contained
in $\KK(S)$ since its elements factor through $\I_C$.  Conversely, take a basis vector $\rho_{R,S}^{U}$ in
\eqref{eq:standard-kernel-admissible-basis}. 
Then its support $U$ satisfies $U\in\cU(S)$ and $U\subsetneq S$. Hence $\rho_{R,S}^{U}= j_{U,S} \circ q_{R,U}$ is contained in the right-hand side. 
This proves
\eqref{eq:standard-kernel-generators}.
\end{proof}




Set
\begin{equation}
\cE_S
=
\{R\in\Int(P)\mid S\in\cD(R)\}.
\label{eq:standard-support-set}
\end{equation}
The following proposition describes the corresponding components of
$\DD(S)$ and the $\Lambda_P$-action explicitly.


\begin{proposition}
\label{prop:standard-concrete-model}
For $R\in \Int(P)$, we have
\begin{equation}
e_R\DD(S)
\cong
\begin{cases}
\kk[q_{R,S}] & \text{if } R\in\cE_S,\\
0            & \text{if } R\notin\cE_S.
\end{cases}
\label{eq:standard-components-explicit}
\end{equation}
Moreover, let $Q,R\in\Int(P)$, let $C\in\Adm(Q,R)$, and assume
$R\in\cE_S$. Then
\begin{equation}
(\rho_{Q,R}^{C})^{\op}\cdot[q_{R,S}]
=
\begin{cases}
[q_{Q,S}] & \text{if } Q\in\cE_S \text{ and } S\subseteq C,\\
0         & \text{otherwise.}
\end{cases}
\label{eq:standard-action-explicit}
\end{equation}
In particular, $\DD(S)$ is a thin $\Lambda_P$-module with support $\cE_S$.
\end{proposition}

\begin{proof}
Fix $R\in\Int(P)$. By
\eqref{eq:projective-injective-component-bases} and
\Cref{lem:standard-kernel-admissible-basis},
\eqref{eq:interval-standard-defining-sequence} gives
\begin{equation}
e_R\DD(S)
\cong
\bigoplus_{U\in\Adm(R,S)}\kk\rho_{R,S}^{U}
\Big/
\bigoplus_{\substack{U\in\Adm(R,S)\\ U\subsetneq S}}
\kk\rho_{R,S}^{U}.
\label{eq:standard-component-basis-quotient}
\end{equation}
Every $U\in\Adm(R,S)$ satisfies $U\subseteq S$. Hence
\eqref{eq:standard-component-basis-quotient} is nonzero precisely when
$S\in\Adm(R,S)$, equivalently $S\in\cD(R)$, equivalently $R\in\cE_S$. 
In this case it is generated by
$[\rho_{R,S}^{S}]=[q_{R,S}]$.
This proves \eqref{eq:standard-components-explicit}.

Let $R\in\cE_S$.  The action of $(\rho_{Q,R}^{C})^{\op}$ on
$[q_{R,S}]$ is represented by $q_{R,S}\circ\rho_{Q,R}^{C}$, and
\Cref{prop:admissible-basis-composition} gives
\begin{equation}
q_{R,S}\circ\rho_{Q,R}^{C}
=
\sum_{W\in\piZero(S\cap C)}\rho_{Q,S}^{W}.
\label{eq:standard-action-before-quotient}
\end{equation}
If $S\subseteq C$, the right-hand side is $q_{Q,S}$; in particular,
$S\in\Adm(Q,S)$ and hence $Q\in\cE_S$.  If $S\nsubseteq C$, every
component of $S\cap C$ is a proper subinterval of $S$ and vanishes in
$\DD(S)$ by \Cref{lem:standard-kernel-admissible-basis}.  This proves
\eqref{eq:standard-action-explicit}.
\end{proof}



\subsubsection{Factorization posets and Loewy structure}
\label{sec:standard-factorization-posets}

We next organize the preceding description in terms of a poset on the support $\cE_S$, which also determines the Loewy structure of $\DD(S)$.

For $R,Q\in\cE_S$, define
\begin{equation}
\label{eq:standard-factorization-order}
R\preceq_S^{\mathcal E}Q
\quad:\Longleftrightarrow\quad
\text{there exists } C\in\Adm(Q,R)
\text{ such that } S\subseteq C.
\end{equation}
For such a component $C$, we have 
$q_{R,S}\circ\rho_{Q,R}^{C}=
q_{Q,S}$.
Hence this relation gives a factorization of $q_{Q,S}$ through $q_{R,S}$.
One checks that $\preceq_S^{\mathcal E}$ is a partial order on
$\cE_S$, with $S$ as its unique minimum.
We call $(\cE_S,\preceq_S^{\mathcal E})$ the
\emph{factorization poset} associated with $\DD(S)$.


The covering relations in this factorization poset correspond
precisely to morphisms between the corresponding interval
representations that are irreducible relative to intervals.

\begin{proposition}
\label{prop:standard-factorization-poset-covers}
Let $R,Q\in\cE_S$ be distinct. Then
\begin{equation}
R\lessdot_S^{\mathcal E}Q
\quad\Longleftrightarrow\quad
\text{there exists a morphism }
\I_Q\longrightarrow\I_R
\text{ irreducible relative to intervals}.
\label{eq:standard-factorization-covers}
\end{equation}
Whenever these conditions hold, such a morphism is unique up to a
nonzero scalar. It is either $q_{Q,R}$ if $R\subsetneq Q$, or
$j_{Q,R}$ if $Q\subsetneq R$.
\end{proposition}

\begin{proof}
Suppose first that
$R\preceq_S^{\mathcal E}Q$, and let
$C\in\Adm(Q,R)$ such that $S\subseteq C$.
Then $C\in\cE_S$, and
$R\preceq_S^{\mathcal E}C
\preceq_S^{\mathcal E}Q$.
If $R\lessdot_S^{\mathcal E}Q$, it follows that $C=R$ or $C=Q$.
Hence $\rho_{Q,R}^{C}$ is respectively $q_{Q,R}$ or $j_{Q,R}$.
The cover condition and
\eqref{eq:standard-factorization-order} show that this morphism is irreducible relative to intervals.

Conversely, let $f\colon\I_Q\to\I_R$ be irreducible relative to intervals.  By \cite[Proposition~6.8]{BBH25}, up to a nonzero scalar, $f$ is either $q_{Q,R}$ or $j_{Q,R}$.  In the first case,
$R\in\Adm(Q,R)$ and $S\subseteq R$, while in the second case,
$Q\in\Adm(Q,R)$ and $S\subseteq Q$.  Hence
$R\prec_S^{\mathcal E}Q$ by
\eqref{eq:standard-factorization-order}. 
The irreducibility of $f$ forces $R\lessdot_S^{\mathcal E}Q$.
\end{proof}

By the description of the Gabriel quiver of $\Lambda_P$ in
\Cref{sec:prelim-intervals},
\Cref{prop:standard-factorization-poset-covers} identifies the Hasse
quiver of $(\cE_S,\preceq_S^{\mathcal E})$ with the full subquiver of
the Gabriel quiver of $\Lambda_P$ on the vertex set $\cE_S$. 

The factorization poset gives a direct diagrammatic description of the
standard module.

\begin{proposition}
\label{prop:standard-factorization-poset-model}
With respect to the bases in \eqref{eq:standard-components-explicit}, 
the module $\DD(S)$ is
represented by the Hasse diagram of
$(\cE_S,\preceq_S^{\mathcal E})$, with a copy of $\kk$ at each vertex
and the identity map along each arrow. All parallel paths have the same
composite.
\end{proposition}


\begin{proof}
By \Cref{prop:standard-concrete-model}, the nonzero components of
$\DD(S)$ are the one-dimensional spaces $\kk[q_{R,S}]$ with
$R\in\cE_S$.  If
$R\lessdot_S^{\mathcal E}Q$, then
\Cref{prop:standard-factorization-poset-covers} and
\eqref{eq:standard-factorization-order} identify the corresponding
arrow with an admissible-component morphism
$\rho_{Q,R}^{C}$ satisfying $S\subseteq C$.
Hence \eqref{eq:standard-action-explicit} sends
$[q_{R,S}]$ to $[q_{Q,S}]$.
After identifying each $[q_{R,S}]$ with $1\in\kk$, every arrow is the
identity map, and the assertion about parallel paths follows.
\end{proof}


Consequently, the faithful quotient
$\Lambda_P/\operatorname{Ann}_{\Lambda_P}\DD(S)$ is isomorphic to the
incidence algebra of $(\cE_S,\preceq_S^{\mathcal E})$. Under this
isomorphism, $\DD(S)$ is isomorphic to the indecomposable projective
module corresponding to the unique minimum $S$ of
$(\cE_S,\preceq_S^{\mathcal E})$.

The description of $\DD(S)$ in terms of the factorization poset also determines its radical filtration. 
For $R\in\cE_S$, let $h_S^{\DD}(R)$ be the largest length of a strict chain from $S$ to $R$ in
$(\cE_S,\preceq_S^{\mathcal E})$.

\begin{corollary}
\label{cor:standard-loewy}
For $n\geq0$,
\begin{equation}
e_R\rad^n\DD(S)
=
\begin{cases}
\kk[q_{R,S}] & \text{if }h_S^{\DD}(R)\geq n,\\
0            & \text{if }h_S^{\DD}(R)<n.
\end{cases}
\label{eq:standard-radical-components}
\end{equation}
Consequently,
\begin{equation}
\rad^n\DD(S)/\rad^{n+1}\DD(S)
\cong
\bigoplus_{\substack{R\in\cE_S\\ h_S^{\DD}(R)=n}}
\LL(R).
\label{eq:standard-radical-layers}
\end{equation}
\end{corollary}



\subsubsection{Morphisms among standard modules}
\label{sec:standard-morphisms}

We now describe the morphisms among standard modules and the resulting
endomorphism algebra of their direct sum.

Suppose that $T\in\cD(S)$. 
The quotient map
$q_{S,T}\colon\I_S\to\I_T$ induces 
$(q_{S,T})_*\colon\PP(S)\to\PP(T)$, see \Cref{sec:prelim-intervals}.

By \Cref{lem:standard-kernel-admissible-basis} and
\eqref{eq:standard-action-explicit}, the composite
\[
\PP(S)
\xrightarrow{(q_{S,T})_*}
\PP(T)
\xrightarrow{\pi_T}
\DD(T)
\]
vanishes on $\KK(S)$ precisely when
no proper relative upset $C\in\cU(S)$ contains $T$, equivalently when $\Min(S)=\Min(T)$.
In this case, it factors uniquely through $\pi_S$ and induces a
nonzero homomorphism
\begin{equation}
\varphi_{S,T}\colon\DD(S)\longrightarrow\DD(T),
\qquad
\varphi_{S,T}([\id_{\I_S}])=[q_{S,T}].
\label{eq:standard-Hom-generator}
\end{equation}

\begin{proposition}
\label{prop:standard-hom-formula}
Let $S,T\in\Int(P)$. Then
\begin{equation}
\Hom_{\Lambda_P}(\DD(S),\DD(T))
=
\begin{cases}
\kk\varphi_{S,T}
& \text{if } T\in\cD(S) \text{ and } \Min(S)=\Min(T),\\
0
& \text{otherwise}.
\end{cases}
\label{eq:standard-hom-formula}
\end{equation}
Moreover, the generator $\varphi_{S,T}$ is a monomorphism in the first case.
\end{proposition}


\begin{proof}
A homomorphism $\DD(S)\to\DD(T)$ is determined by the image of
$[\id_{\I_S}]$, and hence by $e_S\DD(T)$. By
\eqref{eq:standard-components-explicit}, this component is zero unless
$T\in\cD(S)$. If $T\in\cD(S)$, then
$e_S\DD(T)=\kk[q_{S,T}]$, so
$\dim_{\kk}\Hom_{\Lambda_P}(\DD(S),\DD(T))\leq1$.
By the preceding observation, the nonzero homomorphism
$\varphi_{S,T}$ exists precisely when
$\Min(S)=\Min(T)$. This proves
\eqref{eq:standard-hom-formula}.

Suppose that $T\in\cD(S)$ and $\Min(S)=\Min(T)$.
Then $\cE_S\subseteq\cE_T$. In fact, if $R\in\cE_S$, then
$S\in\cD(R)$. Since $T\in\cD(S)$, we have $T\in\cD(R)$, and hence
$R\in\cE_T$.

By \Cref{prop:standard-concrete-model}, for every $R\in\cE_S$ the
components $e_R\DD(S)$ and $e_R\DD(T)$ are generated by
$[q_{R,S}]$ and $[q_{R,T}]$, respectively. Composition of the canonical
quotient maps gives
$e_R\varphi_{S,T}([q_{R,S}])=[q_{R,T}]$.
Thus $e_R\varphi_{S,T}$ is injective for every $R\in\cE_S$, and hence
$\varphi_{S,T}$ is a monomorphism.
\end{proof}

The monomorphism $\varphi_{S,T}$ also has a direct interpretation in
terms of the factorization posets.  Suppose that
$T\in\cD(S)$ and $\Min(S)=\Min(T)$. Then 
$\cE_S\subseteq\cE_T$, and for $R,Q\in\cE_S$ one has
\[
R\preceq_S^{\mathcal E}Q
\quad\Longleftrightarrow\quad
R\preceq_T^{\mathcal E}Q.
\]
Thus $(\cE_S,\preceq_S^{\mathcal E})$ may naturally be regarded as a
full subposet of $(\cE_T,\preceq_T^{\mathcal E})$.  Moreover, by the
definition of the partial order,
\[
\cE_S
=
\{R\in\cE_T\mid S\preceq_T^{\mathcal E}R\},
\]
so this full subposet is the principal upset generated by $S$.  Hence,
by \Cref{prop:standard-factorization-poset-model},
$\varphi_{S,T}$ identifies $\DD(S)$ with the submodule of $\DD(T)$
represented by this principal upset.

We now use the preceding description of morphisms to identify the
endomorphism algebra of the direct sum of the standard modules. For
this purpose, define a relation $\preceq_{\Delta}$ on $\Int(P)$ by the
following equivalent conditions.
\begin{equation}
\begin{aligned}
S\preceq_{\Delta}T
&:\Longleftrightarrow
\Hom_{\Lambda_P}(\DD(S),\DD(T))\neq0\\
&\Longleftrightarrow
T\in\cD(S) \qand \Min(S)=\Min(T)\\
&\Longleftrightarrow
\cE_S\subseteq\cE_T
\qand
\cE_S\text{ is an upset in }
(\cE_T,\preceq_T^{\cE}).
\end{aligned}
\label{eq:standard-morphism-order-support}
\end{equation}
The relation $\preceq_{\Delta}$ is a partial order on $\Int(P)$.

\begin{corollary}
\label{cor:standard-endomorphism-incidence}
The endomorphism algebra
\[
\End_{\Lambda_P}
\left(
\bigoplus_{S\in\Int(P)}\DD(S)
\right)
\]
is isomorphic to the incidence algebra of
$(\Int(P),\preceq_{\Delta})$.
\end{corollary}

\begin{proof}
By \Cref{prop:standard-hom-formula}, the morphisms
$\varphi_{S,T}$ with $S\preceq_{\Delta}T$ form a basis of the
endomorphism algebra. Whenever
$S\preceq_{\Delta}T\preceq_{\Delta}U$, their definitions give
\[
\varphi_{T,U}\circ\varphi_{S,T}
=
\varphi_{S,U}.
\]
Thus their composition agrees with the multiplication of the incidence
basis of $(\Int(P),\preceq_{\Delta})$.
\end{proof}

\subsection{Costandard modules via poset duality}
\label{sec:costandard-via-duality}

We record the corresponding description of the costandard modules.
The poset duality \eqref{eq:poset-duality}, together with
\eqref{eq:QH-opposite-standard-costandard-duality}, gives
\begin{equation}
\NN_P(S)
\cong
\D\DD_{P^{\op}}(S).
\label{eq:interval-standard-costandard-duality}
\end{equation}

For $S\in\Int(P)$, set
\begin{equation}
\cM_S
=
\{R\in\Int(P)\mid S\in\cU(R)\}.
\label{eq:costandard-support-set}
\end{equation}

\begin{proposition}
\label{prop:costandard-concrete-model}
Let $R,S\in\Int(P)$. Then
\begin{equation}
e_R\NN(S)
\cong
\begin{cases}
\kk j_{S,R}^{\vee} & \text{if } R\in\cM_S,\\
0                  & \text{if } R\notin\cM_S.
\end{cases}
\label{eq:costandard-components-explicit}
\end{equation}
Moreover, let $Q,R\in\Int(P)$, let $C\in\Adm(Q,R)$, and assume
$R\in\cM_S$. Then
\begin{equation}
(\rho_{Q,R}^{C})^{\op}\cdot j_{S,R}^{\vee}
=
\begin{cases}
j_{S,Q}^{\vee}
& \text{if } Q\in\cM_S \text{ and } S\subseteq C,\\
0
& \text{otherwise}.
\end{cases}
\label{eq:costandard-action-explicit}
\end{equation}
In particular, $\NN(S)$ is thin, with support $\cM_S$.
\end{proposition}

\begin{proof}
In $P^{\op}$, relative downsets are relative upsets in $P$, so the
standard support set for $S$ in $P^{\op}$ is $\cM_S$. Pointwise
duality sends the quotient map $q_{R,S}$ in $\rep_{\kk}P^{\op}$ to
the inclusion map $j_{S,R}$ in $\rep_{\kk}P$. Hence
\eqref{eq:costandard-components-explicit} follows from
\eqref{eq:interval-standard-costandard-duality} and
\eqref{eq:standard-components-explicit} applied to $P^{\op}$.

For the action, regard $\NN(S)$ as a submodule of
$\II(S)=\D(e_S\Lambda_P)$. If $Q\notin\cM_S$, then
$e_Q\NN(S)=0$. Suppose that $Q\in\cM_S$. The coefficient of
$j_{S,Q}^{\vee}$ in
$(\rho_{Q,R}^{C})^{\op}\cdot j_{S,R}^{\vee}$ is obtained by
evaluating $j_{S,R}^{\vee}$ on
$\rho_{Q,R}^{C}\circ j_{S,Q}$. By
\Cref{prop:admissible-basis-composition}, the coefficient of
$j_{S,R}=\rho_{S,R}^{S}$ in this composite is one precisely when
$S\subseteq C$. This proves \eqref{eq:costandard-action-explicit}.
\end{proof}

For $R,Q\in\cM_S$, define
\begin{equation}
\label{eq:costandard-factorization-order}
R\preceq_S^{\mathcal M}Q
\quad:\Longleftrightarrow\quad
\text{there exists } C\in\Adm(Q,R)
\text{ such that } S\subseteq C.
\end{equation}
For such a component $C$, we have
$\rho_{Q,R}^{C}\circ j_{S,Q}=j_{S,R}$.
This is a partial order on $\cM_S$, and $S$ is its unique maximum.
Under \eqref{eq:poset-duality}, it is the order dual of the
factorization poset for the standard module $\DD_{P^{\op}}(S)$.

Applying the results of
\Cref{sec:standard-factorization-posets} to $P^{\op}$ and dualizing,
the module $\NN(S)$ is represented by the Hasse diagram of
$(\cM_S,\preceq_S^{\mathcal M})$, with a copy of $\kk$ at each
vertex and the identity map along each arrow; all parallel paths have
the same composite. Consequently,
$\Lambda_P/\operatorname{Ann}_{\Lambda_P}\NN(S)$ is isomorphic to
the incidence algebra of $(\cM_S,\preceq_S^{\mathcal M})$, and
$\NN(S)$ is isomorphic to the indecomposable injective module
corresponding to its unique maximum $S$.

The same duality determines the socle filtration of $\NN(S)$ from
the lengths of chains ending at $S$. It also gives the dual statements
for morphisms among costandard modules from
\Cref{sec:standard-morphisms}; in particular, every nonzero morphism
between costandard modules is an epimorphism.


\subsection{Projective and injective filtrations and bimodule strata}
\label{sec:standard-cardinality-filtrations}

In this subsection, we use the heredity chain
\eqref{eq:interval-cardinality-heredity-chain} to construct standard
and costandard filtrations of the indecomposable projective and
injective modules, respectively, and to identify the successive strata
$J_t/J_{t-1}$ as bimodules.


Fix $S\in\Int(P)$ and set $n=|S|$.
For $0\leq t\leq n$, set
\[
F_t(S)
:=
J_t e_S
=
\operatorname{span}_{\kk}
\left\{
\rho_{R,S}^{C}
\;\middle|\;
R\in\Int(P),\ C\in\Adm(R,S),\ |C|\leq t
\right\}
\subseteq\PP(S).
\]
Since $J_t$ is a two-sided ideal, $F_t(S)$ is a left
$\Lambda_P$-submodule of $\PP(S)$.

Dually, for $1\leq t\leq n+1$, set
\[
G_t(S)
:=
\D(e_S\Lambda_P/e_SJ_{t-1})
=
\operatorname{span}_{\kk}
\left\{
(\rho_{S,R}^{C})^{\vee}
\;\middle|\;
R\in\Int(P),\ C\in\Adm(S,R),\ |C|\geq t
\right\}
\subseteq\II(S),
\]
where $\D(e_S\Lambda_P/e_SJ_{t-1})$ is viewed naturally as a
submodule of $\II(S)=\D(e_S\Lambda_P)$.

\begin{proposition}
\label{prop:projective-injective-cardinality-filtrations}
The left $\Lambda_P$-submodules defined above give filtrations
\begin{equation}
0=F_0(S)\subseteq\cdots\subseteq F_{n-1}(S)=\KK(S)
\subseteq F_n(S)=\PP(S)
\label{eq:projective-cardinality-filtration-sequence}
\end{equation}
and
\begin{equation}
0=G_{n+1}(S)\subseteq G_n(S)=\NN(S)
\subseteq\cdots\subseteq G_1(S)=\II(S).
\label{eq:injective-cardinality-filtration-sequence}
\end{equation}
Their successive factors are
\begin{align}
F_t(S)/F_{t-1}(S)
&\cong
\bigoplus_{\substack{C\in\cU(S)\\ |C|=t}}\DD(C),
\label{eq:projective-standard-factors}\\
G_t(S)/G_{t+1}(S)
&\cong
\bigoplus_{\substack{C\in\cD(S)\\ |C|=t}}\NN(C).
\label{eq:injective-costandard-factors}
\end{align}
In particular, all standard and costandard filtration multiplicities in
these filtrations are zero or one.
\end{proposition}

\begin{proof}
Since $J_\bullet$ is an increasing chain of two-sided ideals, the
modules $F_t(S)=J_te_S$ form an increasing filtration of $\PP(S)$.
The admissible-component basis gives $F_0(S)=0$ and
$F_n(S)=\PP(S)$, while
\eqref{eq:standard-kernel-admissible-basis} gives
$F_{n-1}(S)=\mathsf K(S)$.

Fix $1\leq t\leq n$. The quotient $F_t(S)/F_{t-1}(S)$ has basis
given by the classes of $\rho_{R,S}^C$ with
$C\in\Adm(R,S)$ and $|C|=t$. 
We show that the contribution with fixed support $C$ is a copy of
$\DD(C)$.

Let $C\in\cU(S)$ with $|C|=t$. Composition with $j_{C,S}$ induces
a map $\PP(C)\rightarrow F_t(S)/F_{t-1}(S)$.
For $R\in\Int(P)$ and $U\in\Adm(R,C)$, the basis vector
$\rho_{R,C}^U$ is sent to the class of $\rho_{R,S}^U$.
If $U\subsetneq C$, then
$|U|<t$, so this class vanishes in the quotient. Hence the map
factors through $\DD(C)$ and gives
$f_C\colon\DD(C)\rightarrow F_t(S)/F_{t-1}(S)$.

We now compare these maps for all $C\in\cU(S)$ with $|C|=t$.
Fix $R\in\Int(P)$. By \Cref{prop:standard-concrete-model},
$e_R\DD(C)$ is one-dimensional, with basis $[q_{R,C}]$, precisely
when $C\in\cD(R)$. Since $C\in\cU(S)$, this is equivalent to
$C\in\Adm(R,S)$, and $f_C$ sends $[q_{R,C}]$ to
$[\rho_{R,S}^C]$.
On the other hand, the classes $[\rho_{R,S}^C]$ with
$C\in\Adm(R,S)$ and $|C|=t$ form a basis of
$e_R(F_t(S)/F_{t-1}(S))$. Thus, for each $R$, the nonzero
components $e_R\DD(C)$, as $C$ varies, are mapped bijectively onto
this basis. This proves the asserted decomposition.

The assertions for the injective filtration follow by applying the
projective statement to $P^{\op}$ and then applying $\kk$-duality.
Under order reversal, relative upsets correspond to relative
downsets, and the resulting filtration is precisely $G_\bullet(S)$.
\end{proof}



We next identify the successive quotients $J_t/J_{t-1}$ as bimodules.
For each $t$, the quotient $J_t/J_{t-1}$ is a heredity ideal of
$\Lambda_P/J_{t-1}$, so the multiplication isomorphism for heredity
ideals of Dlab--Ringel \cite[Proposition~7]{DRconstruction89} applies.
This gives the following decomposition in terms of standard and costandard modules.


\begin{proposition}
\label{prop:cardinality-bimodule-strata}
For $1\leq t\leq |P|$, multiplication induces an isomorphism of
$\Lambda_P$--$\Lambda_P$-bimodules
\begin{equation}
\label{eq:cardinality-stratum-vector-space-decomposition}
\theta_t\colon
\bigoplus_{\substack{T\in\Int(P)\\ |T|=t}}
\DD(T)\otimes_\kk\D\NN(T)
\xrightarrow{\sim}
J_t/J_{t-1}.
\end{equation}
Moreover, for each $T\in\Int(P)$ with $|T|=t$, the image under
$\theta_t$ of the summand
$\DD(T)\otimes_\kk\D\NN(T)$ is
\[
\operatorname{span}_{\kk}
\left\{
(\rho_{R,S}^{T})^{\op}+J_{t-1}
\;\middle|\;
R,S\in\Int(P),\ T\in\Adm(R,S)
\right\}.
\]
\end{proposition}

\begin{proof}
Fix $t$, set $A_t=\Lambda_P/J_{t-1}$, and write $e_T$ and
$\varepsilon_t$ also for their images in $A_t$.  By
\eqref{eq:interval-cardinality-heredity-chain}, $A_t$ is
quasi-hereditary and
\[
A_t\varepsilon_tA_t=J_t/J_{t-1}
\]
is a heredity ideal of $A_t$.  Hence
\cite[Proposition~7]{DRconstruction89} gives an isomorphism of
$A_t$--$A_t$-bimodules
\begin{equation}
\label{eq:cardinality-heredity-multiplication}
A_t\varepsilon_t
\otimes_{\varepsilon_tA_t\varepsilon_t}
\varepsilon_tA_t
\xrightarrow{\sim}
J_t/J_{t-1}
\end{equation}
induced by multiplication.

By \eqref{eq:interval-cardinality-admissible-basis-layer},
\[
\varepsilon_tA_t\varepsilon_t
=
\bigoplus_{\substack{T\in\Int(P)\\ |T|=t}}\kk e_T.
\]
Therefore the source of the multiplication map \eqref{eq:cardinality-heredity-multiplication} decomposes as
\[
\bigoplus_{\substack{T\in\Int(P)\\ |T|=t}}
A_te_T\otimes_\kk e_TA_t.
\]
For $|T|=t$, \Cref{prop:projective-injective-cardinality-filtrations} gives $F_{t-1}(T) = \KK(T)$ and $G_t(T) = \NN(T)$. Hence 
\[
A_te_T
\cong
\PP(T)/\KK(T)
=
\DD(T),
\qquad
e_TA_t
\cong
\D\NN(T).
\]
Under these identifications, the multiplication map is
\eqref{eq:cardinality-stratum-vector-space-decomposition}.

Using the canonical double-dual identification, write $j_{T,S}$ for
the basis vector of $\D\NN(T)$ dual to $j_{T,S}^{\vee}$.
For $T\in\Adm(R,S)$, multiplication gives
\begin{equation}
\label{eq:cardinality-bimodule-basis-correspondence}
\theta_t\bigl([q_{R,T}]\otimes j_{T,S}\bigr)
=
q_{R,T}^{\op}j_{T,S}^{\op}+J_{t-1}
=
(\rho_{R,S}^{T})^{\op}+J_{t-1}.
\end{equation}
By \eqref{eq:interval-cardinality-admissible-basis-layer}, these are
precisely the classes of the admissible-component basis elements with
support $T$, which proves the second assertion.
\end{proof}


